% !PRINTED 288
\weberpage{288}
%<<CONT
sein, und wir wollen ausserdem noch annehmen, dass sie durch
keinen der Schnittpunkte von $\varphi$ und $\psi$ hindurchgeht.

Wir geben diesen drei Functionen eine bestimmte cyklische
Reihenfolge
\[
f,\ \varphi,\ \psi,
\]
so dass auf $\psi$ wieder $f$ folgen soll, und verbinden mit dieser
Reihenfolge eine bestimmte, etwa die positive Fortschrittsrichtung
auf jeder der drei Curven. Wir wollen festsetzen, dass mit der
umgekehrten Reihenfolge
\[
f,\ \psi,\ \varphi
\]
auf jeder der drei Curven die entgegengesetzte, also die negative,
Fortschrittsrichtung verbunden sei.

Wir durchlaufen nun, bei der ersten Reihenfolge, die Curve
$f$ in positivem Sinne, und achten auf die Schnittpunkte von $f$
und $\varphi$. Ein solcher soll ein Austrittspunkt $A(f; \varphi, \psi)$ genannt

\weberfigureimage{weber-b1-fig08.png}{0.48}

werden, wenn die positive Richtung von $f$ aus dem Inneren des
Binnenraumes $(\varphi, \psi)$ in das Aeussere, und ein Eintrittspunkt
$E(f; \varphi, \psi)$, wenn er aus dem Aeusseren in das Innere führt.

In der Fig.~8 sind die Punkte $a$ die Austrittspunkte, der Punkt
$e$ ein Eintrittspunkt.

Die Gesammtanzahl der Ein- und Austrittspunkte stimmt überein
mit der Anzahl der Schnittpunkte von $f$ und $\varphi$, und ist also, da
beide Curven geschlossen sind, eine \gesperrt{gerade Zahl}. Ist $a$ die Anzahl
der Punkte $A(f; \varphi, \psi)$, $e$ die Anzahl der Punkte $E(f; \varphi, \psi)$, so
ist also auch $e - a$ eine gerade Zahl und
\[
\tfrac{1}{2}(e - a) = k
\]
eine ganze Zahl, die nach \gesperrt{Kronecker} die \gesperrt{Charakteristik
des Functionensystems} $f$, $\varphi$, $\psi$ heisst. (Im Falle der Fig.~8
ist sie gleich $-1$).

Wenn wir die Functionen $f$, $\varphi$, $\psi$ cyklisch vertauschen, so
erhalten wir drei Bestimmungen für die Charakteristik, und wenn
wir die Reihenfolge umkehren und dabei nach der getroffenen
Vereinbarung auch die positiven Richtungen durch die negativen
ersetzen, drei weitere. Diese Bestimmungen sind, wenn wir jetzt
die Symbole $A(f; \varphi, \psi)$, $E(f; \varphi, \psi)$ zugleich als Bezeichnung
für die Anzahlen der betreffenden Punkte brauchen,
%CONT>>
